\chapter{Esempio di codice sviluppato}
\label{cap:appendice}

In questa Appendice sono riportati i principali frammenti di codice utilizzati per simulare e rappresentare i sistemi in Figura \ref{fig:SIE_sim}. Come nel resto del progetto, il codice e i relativi commenti sono in lingua inglese.

In primo luogo, si è istanziata la sorgente imperturbata, utilizzando i parametri e i profili definiti nelle Tabelle \ref{tab:lens_par_sim} e \ref{tab:source_par_sim} e si sono salvate in variabili globali le relative immagini di interesse.
Successivamente, si è definita la funzione \texttt{make\_composite\_lens}, che costruisce sia il satellite sulla base del profilo di densità di massa scelto dall'utente (SIE o NFW), sia la lente composita. Si sono poi definiti la posizione del perturbatore, in quanto mantenuta costante nei due casi, e i parametri di base dei due profili.
\begin{lstlisting}
    source_unperturbed = sersic(size=fov, Npix=npix, gl=lens, save_unlensed=True, **source_kwargs)
    total_image_unperturbed = source_unperturbed.image + lens_light.image
    difference_unperturbed = source_unperturbed.image - source_unperturbed.image

    # Create a composite lens made of the main lens plus either a SIE or NFW perturber
    def make_composite_lens(
    base_lens,
    perturbation_kind,
    perturbation_kwargs,
    co,
    zl,
    zs,
    theta,
    ):
        if perturbation_kind == "SIE":
            perturbation = sie(
            co=co,
            zl=zl,
            zs=zs,
            **perturbation_kwargs 
            )

        if perturbation_kind == "NFW":
            perturbation = nfwell(
          co=co,
          zl=zl,
          zs=zs,
          **perturbation_kwargs 
        )
        perturbation.setGrid(theta=theta)

        composite = compositeModel(
            co,
            models=[base_lens, perturbation]
        )

        composite.zl = zl
        composite.zs = zs
        composite.dl = base_lens.dl
        composite.ds = base_lens.ds
        composite.dls = base_lens.dls
        composite.setGrid(theta=theta)

        return perturbation, composite

    pert_pos = [1, -1.05]

    second_sie_kwargs = {
    "sigma0": 90.0,
    "q": 0.9,
    "pa": np.deg2rad(-20.0),
    "theta_c": 1.0e-4,
    "x1": pert_pos[0], 
    "x2": pert_pos[1], 
    }

    second_nfw_kwargs = {
    "mass": 5.0e11, # total mass
    "conc": 11.0, # concentration
    "q": 0.95,
    "pa": np.deg2rad(-20.0),
    "x1": pert_pos[0],
    "x2": pert_pos[1],
    }
\end{lstlisting}
La funzione appena definita viene quindi utilizzata per aggiungere alla lente principale un satellite SIE con dispersione di velocità variabile, salvando le lenti composite e le relative linee critiche per poterle successivamente rappresentare.
\begin{lstlisting}
    sigma0_values_pert = [50.0, 70.0, 90.0]
    second_lens_SIE_all = []
    composite_lens_SIE_all = []
    composite_critical_lines_SIE = []
    
    for sigma0_pert in sigma0_values_pert:
        second_sie_kwargs_var = {
        "sigma0": sigma0_pert, 
        "q": 0.9,
        "pa": np.deg2rad(-20.0),
        "theta_c": 1.0e-4,
        "x1": pert_pos[0], 
        "x2": pert_pos[1], 
        }
    
        second_lens_SIE, lens_with_second_SIE = make_composite_lens(lens, "SIE", second_sie_kwargs_var, co, zl, zs, theta)

        second_lens_SIE_all.append(second_lens_SIE)
        composite_lens_SIE_all.append(lens_with_second_SIE)

        critical_lines_with_second_SIE = [lens_with_second_SIE.getCritPoints(cl) for cl in lens_with_second_SIE.tancl()]
        composite_critical_lines_SIE.append(critical_lines_with_second_SIE)
\end{lstlisting}
Il blocco di codice seguente esegue il ray tracing attraverso la lente composita e salva le immagini simulate, le mappe di convergenza e i residuali in apposite liste, che verranno successivamente utilizzate per la rappresentazione grafica.
\begin{lstlisting}
    composite_sources_all_SIE = []
    composite_total_images_all_SIE = []
    composite_kappa_maps_all_SIE = []
    SIE_residuals = []
    
    for i, pert in enumerate(second_lens_SIE_all):
        source_with_second_SIE = sersic(
        size=fov,
        Npix=npix,
        gl=composite_lens_SIE_all[i],
        save_unlensed=True,
        **source_kwargs,
        )
        composite_sources_all_SIE.append(source_with_second_SIE)
        composite_kappa_maps_all_SIE.append(composite_lens_SIE_all[i].ka)
    
        total_image_with_second_SIE = source_with_second_SIE.image + lens_light.image
        composite_total_images_all_SIE.append(total_image_with_second_SIE)
    
        difference_SIE = source_with_second_SIE.image - source_unperturbed.image
        SIE_residuals.append(difference_SIE)
        
\end{lstlisting}
Infine, questo implementa la visualizzazione grafica del risultato:
\begin{lstlisting}
    vmax_image_SIE = percentile_vmax(composite_total_images_all_SIE)
    vmax_kappa_SIE = percentile_vmax(composite_kappa_maps_all_SIE)
    vmax_difference_SIE = percentile_vmax(SIE_residuals)
    
    fig3, ax3 = plt.subplots(3, 4, figsize=(15, 10), constrained_layout=True)
    
    # Effect of varying the velocity dispersion of a SIE perturbator
    fig3.suptitle("Effetto della variazione della dispersione di velocita di un perturbatore SIE", fontsize=20)
    
    for ax in ax3.flat:
        ax.set_label("arcsec")
        ax.set_aspect("equal")
    
    # Unperturbed references
    ax3[0,0].set_title("Immagine imperturbata")
    ax3[0,0].imshow(total_image_unperturbed, origin="lower", extent=extent, cmap="magma", vmax=vmax_image_SIE)
    ax3[1,0].set_title('Convergenza $\kappa$ imperturbata')
    ax3[1,0].imshow(lens.ka, origin="lower", extent=extent, cmap="coolwarm", vmax=vmax_kappa_SIE)
    ax3[2,0].set_title("Differenza puntuale - imperturbata")
    ax3[2,0].imshow(difference_unperturbed, origin="lower", extent=extent, cmap="PuOr", vmin=-vmax_difference_SIE, vmax=vmax_difference_SIE)
    
       
    for xcritu, ycritu in critical_lines_lens:
        ax3[0,0].plot(xcritu, ycritu, color="white", linestyle="dotted")
        ax3[1,0].plot(xcritu, ycritu, color="white")
        ax3[2,0].plot(xcritu, ycritu, color="indigo", linestyle="dotted")
    
    for i, sigma0_pert in enumerate(sigma0_values_pert):
        # Perturbed images and convergence maps 
        ax3[0,i+1].set_title(
            rf'$\sigma_\mathrm{{pert}}={sigma0_pert}\,\mathrm{{km\,s^{{-1}}}}$, '
            rf'$\theta_\mathrm{{E, pert}}={second_lens_SIE_all[i].bsie():.2f}$"')
        im0 = ax3[0, i+1].imshow(composite_total_images_all_SIE[i], origin="lower", extent=extent, cmap="magma", vmax=vmax_image_SIE)
    
        ax3[1,i+1].set_title(
            rf'Convergenza $\kappa$, $\sigma_\mathrm{{pert}}={sigma0_pert}\,\mathrm{{km\,s^{{-1}}}}$')
        im1 = ax3[1,i+1].imshow(composite_kappa_maps_all_SIE[i], origin="lower", extent=extent, cmap="coolwarm", vmax=vmax_kappa_SIE)
    
        ax3[2,i+1].set_title(
            rf'Differenza puntuale, $\sigma_\mathrm{{pert}}={sigma0_pert}\,\mathrm{{km\,s^{{-1}}}}$'
        )
        im2 = ax3[2, i+1].imshow(SIE_residuals[i], origin = "lower", extent=extent, cmap="PuOr", vmin=-vmax_difference_SIE, vmax=vmax_difference_SIE)
    
        for xcrit, ycrit in composite_critical_lines_SIE[i]:
            ax3[0,i+1].plot(xcrit, ycrit, color="white", linestyle="dotted")
            ax3[1,i+1].plot(xcrit, ycrit, color="white")
            ax3[2,i+1].plot(xcrit, ycrit, color="indigo", linestyle="dotted")
    
    # Surface brightness
    cbar0 = fig3.colorbar(
        im0,
        ax=ax3[0, :],
        shrink=0.85,
        pad=0.02
    )
    cbar0.set_label("Luminosita superficiale")
    
    # Convergence
    cbar1 = fig3.colorbar(
        im1,
        ax=ax3[1, :],
        shrink=0.85,
        pad=0.02
    )
    cbar1.set_label(r"Convergenza $\kappa$")
    
    # Surface brightness (again)
    cbar2 = fig3.colorbar(
        im2,
        ax=ax3[2, :],
        shrink=0.85,
        pad=0.02
    )
    cbar2.set_label("Luminosita superficiale")
\end{lstlisting}
Per uniformare la scala di visualizzazione delle immagini appartenenti alla stessa riga, era stata  definita in precedenza la funzione
\begin{lstlisting}
    def percentile_vmax(images, percentile=99.5):
    return max(
        np.percentile(img[np.isfinite(img)], percentile)
        for img in images
    )
\end{lstlisting}
